\documentclass[11pt]{article}
\usepackage{mathblatt}
\makeatletter
\setlength{\mbpfbe}{0mm}% keine Punkte (Bauregel 6.4)
\makeatother
% Rechenraum als Karo (Bauregel 6.9): n Rechenschritte = n mal 10 mm
\newcommand{\karo}[1]{\par\smallskip\noindent\begin{tikzpicture}
  \clip (0,0) rectangle ({\linewidth-1mm},{#1*10mm});
  \draw[black!16,line width=0.3pt,step=5mm] (0,0) grid ({\linewidth},{#1*10mm});
  \end{tikzpicture}\par}
\newcommand{\linie}[1]{\,{\color{black!45}\rule[-1pt]{#1}{0.4pt}}\,}
\newcommand{\pkt}{\textperiodcentered{}}
\begin{document}

\pfheftstil
\pflaufkopf[P10 \pkt{} mehr]{Parabel: Scheitelpunkt und Skizze \pkt{} Lösungen}{PA30-4}\pfseitevon
\pfheftkopf{Parabel: Scheitelpunkt und Skizze}{Lösungen \pkt{} mehr}
\begin{pfloesungb}{}
\lzb{1.}{}{S($-$3|1)}{S($-$3|1)}
\end{pfloesungb}
\begin{pfloesungb}{}
\lzb{2.}{}{S($-$3|$-$5)}{S($-$3|$-$5)}
\end{pfloesungb}
\begin{pfloesungb}{}
\lzb{3.}{}{S($-$1|$-$3)}{S($-$1|$-$3)}
\end{pfloesungb}
\begin{pfloesungb}{}
\lzb{4.}{}{S($-$4|$-$3)}{S($-$4|$-$3)}
\end{pfloesungb}
\begin{pfloesungb}{}
\lzb{5.}{}{y = (x $-$ 4)² + 2}{y = (x $-$ 4)² + 2}
\end{pfloesungb}
\begin{pfloesungb}{}
\lzb{6.}{}{S(3|0)}{S(3|0)}
\end{pfloesungb}
\begin{pfloesungb}{}
\lzb{7.}{}{p(x) = (x $-$ 3)² $-$ 5}{p(x) = (x $-$ 3)² $-$ 5}
\end{pfloesungb}
\begin{pfloesungb}{}
\lzb{8.}{}{p(x) = (x + 1)² $-$ 3}{p(x) = (x + 1)² $-$ 3}
\end{pfloesungb}
\begin{pfloesungb}{}
\lzb{9.}{}{q(x) = $-$(x + 4)²}{q(x) = $-$(x + 4)²}
\end{pfloesungb}
\begin{pfloesungb}{}
\lzb{10.}{}{h(t) = $-$2t² + 2750 (Start bei 2750 m, Höhe sinkt: Parabel nach unten geöffnet)}{h(t) = $-$2t² + 2750 (Start bei 2750 m, Höhe sinkt: Parabel nach unten geöffnet)}
\end{pfloesungb}
\begin{pfloesungb}{}
\lzb{11.}{}{Beispiel: f(x) = 1 (jede waagerechte Gerade oberhalb des Scheitelpunkts trifft p nicht)}{Beispiel: f(x) = 1 (jede waagerechte Gerade oberhalb des Scheitelpunkts trifft p nicht)}
\end{pfloesungb}
\begin{pfloesungb}{}
\lzb{12.}{}{q$'$(x) = $-$x² $-$ 2}{Beispiel: q(x) = x² + 2; q$'$(x) = $-$x² $-$ 2}
\end{pfloesungb}
\end{document}